\begin{abstract}
Multi-tenant retrieval-augmented generation (RAG) applications copy access-control state and
document content into derived artifacts: ACL copies in the vector index, semantic-cache entries,
and chat memory with model-written summaries. Revoking a user's access in the identity system does
not revoke these copies. We measure what they disclose after revocation, separating secrets the user
had already been shown (re-disclosure, a violation of complete mediation) from secrets the user had
never seen, including content of a document that changed after the revocation. With real models
from \pending{3B to 32B} on generated and real (Enron) corpora, \pending{first result sentence:
new disclosure under baseline and through which path (stale index, shared cache), written from
p4\_defenses and p4\_shared\_state}; \pending{memory lifetime under session retention, from
p2\_lifetimes}. Two published retrieval-time enforcement designs remove the index channel
and leave \pending{Y} of the leakage. We then compare read-time defenses for derived state:
provenance taint, a per-user permission epoch, and session flush. \pending{defense result: which
reach zero leakage, including with lost revocation hooks and document reclassification, and what
each costs the revoked user and bystanders; name the summary design (model-written or extractive)
for any sentence-level-taint claim}. \pending{existence result from p3\_existence, insider and
outsider}. A user who pastes back what they once received defeats every
artifact-level defense, and we measure how far that reaches.
\end{abstract}
